\section{理解与生成能否正迁移}

前面展示了许多“统一模型”，但 architecture unification 不等于 capability transfer。一个模型可能使用同一 backbone 完成 VQA 与 image generation，却仍需要两套 representation；也可能生成逼真视频，却没有更好的 spatial reasoning。讲座后半段把这种不确定性放到中心，并提出语言与视觉信号的信息结构差异可能是根因之一。

\subsection{先把 transfer direction 画清楚}

前面所有“统一模型”都可能让人默认任务会互相帮助。本节先纠正这个直觉：transfer 有方向，也需要 matched single-task baseline；slide 047--048 因此负责定义实验，而不是宣布答案。

\lecturefigure{slide-047.jpg}{第二条研究主线：understanding 与 generation 是否正迁移}{官方 deck 第 47 页；课堂 00:35:11--00:35:30}

\begin{importantbox}{Positive transfer 的可检验定义}
在控制数据量、compute、参数和评价协议后，联合训练使目标任务 $B$ 的性能优于只训练 $B$ 的 baseline，才可称任务 $A\rightarrow B$ positive transfer。若只是总模型更大、训练更久或使用额外标注，不能把增益归因于 transfer。
\end{importantbox}

\lecturefigure{slide-048.jpg}{Understanding $\rightarrow$ Generation 与 Generation $\rightarrow$ Understanding 的两条路径}{官方 deck 第 48 页；课堂 00:35:30--00:36:10}

\begin{knowledgebox}{读图：两条箭头需要不同实验}
左侧 understanding-to-generation 假设 perception、captioning 或 reasoning representation 能提高 prompt fidelity、layout 与 editing；右侧 generation-to-understanding 假设生成式 world modeling 帮助模型掌握对象、运动或因果。要验证左箭头，需要固定 generator 并加入理解训练；验证右箭头，则要固定理解 benchmark 并加入生成目标。不能用同一组 qualitative samples 同时证明两条方向。
\end{knowledgebox}

\begin{warningbox}{共享参数不是 transfer 的证据}
两个任务经过同一 Transformer 只说明它们可能交互。若 gradient 冲突、loss scale 不平衡或 representation 不兼容，联合训练甚至可能 negative transfer。MoT 的结果已提供反例边界：image generation specialization 有收益，但没有观察到相应 image-understanding 提升。
\end{warningbox}

\subsection{为什么语言更像“压缩后的认知接口”}

语言 token 往往已经经过人类选择：一句话保留对象、关系、意图与异常，丢掉大量像素细节。图像和视频则接近 raw sensory observation，包含照明、纹理、背景、相机运动和重复帧。于是同样的 next-step prediction 在两类数据上拥有不同 signal-to-noise ratio。

\lecturefigure{slide-051.jpg}{语言压缩、视觉被动观察、loss geometry 与帧冗余}{官方 deck 第 51 页；课堂 00:36:10--00:39:11}

\begin{knowledgebox}{读图：四个气泡构成一个因果假说，而非已证定律}
左侧指出 language 是高度压缩的 cognition abstraction；右侧列出 image/video 的三个困难：被动观察、loss landscape 与 frame correlation。它们共同解释为什么视觉 loss 可能很低但生成仍不自然，也解释为什么视频拥有海量 bytes 却未必提供等量独立 supervision。图支持研究假设，尚不能证明语言天然优于一切视觉表示。
\end{knowledgebox}

若相邻视频 frame 的相关系数近似为 $\rho$，一个常用的 AR(1) 直觉是
$$
N_{\text{eff}}\approx N\frac{1-\rho}{1+\rho}.
$$
$N$ 是 frame 数，$N_{\text{eff}}$ 是近似独立样本数。$\rho$ 越接近 1，新增 frame 带来的独立信息越少。这个公式只是统计直觉，不适合直接用于复杂视频；它提醒我们“更多 frames”不等于“更多 semantic events”。

\begin{warningbox}{低 perceptual loss 不保证高人类质量}
Pixel/latent MSE 会平均化多个可能未来，diffusion loss 又在噪声空间优化。某些误差对人类视觉极其显著，却只占很小数值；另一些像素变化数值大但语义不变。因此必须结合 human preference、object relation、temporal consistency 与 task success，而不能只追最低 training loss。
\end{warningbox}

\teachervoice{00:38:11--00:38:21，老师指出 image/video generation 即使在人看来还不好，loss 也可能已经开始变好。这个提醒把 optimization metric 与 user-visible quality 明确分开。}

\subsection{Digital multimodal processing 仍不是 physical intelligence}

前一节解释了视觉学习信号的困难，本节把范围从 optimization 扩展到 system capability。slide 052 要回答的是：即使模型能处理和生成多种数字内容，距离持续感知、控制与行动还缺哪些层。

\lecturefigure{slide-052.jpg}{从数字信息处理到更广泛的物理世界多模态智能}{官方 deck 第 52 页；课堂 00:39:12--00:40:04}

\begin{knowledgebox}{读图：中间模型只覆盖整张图的一部分}
左侧是 text、image、video、audio 的 digital information；中间列出 Chameleon、Transfusion、MoT 与 cross-modal generation；右侧 physical world 还要求 spatial awareness、real-time world understanding、control、exploration 与 long-horizon interaction。底部结论很克制：today's models excel at multimodal information processing, but fall far short of full physical-world intelligence。
\end{knowledgebox}

物理闭环可以写成
$$
o_t\xrightarrow{\text{encode}}h_t
\xrightarrow{\text{reason/plan}}a_t
\xrightarrow{\text{environment}}o_{t+1}.
$$
数字生成通常只评价 $a_t$（输出内容）；机器人系统还要评价环境转移是否安全、状态估计是否漂移、延迟是否满足控制周期，以及失败后能否 recovery。

\begin{importantbox}{从 omni model 到 embodied intelligence 还缺什么}
缺少的不只是更多 modality，而是 active perception、state estimation、spatial-temporal memory、closed-loop control、uncertainty calibration、real-time infrastructure 与 safety constraints。把一张图送进 LLM 和让机器人持续看、听、动，系统要求相差一个完整闭环。
\end{importantbox}

\subsection{本章小结}

Understanding 与 generation 的 transfer 必须按方向、baseline 和 matched compute 验证。语言提供高密度抽象信号，视觉提供高带宽但高冗余的感知数据；二者的最佳 representation 和 loss 未必相同。当前 omni model 在数字信息处理中进步显著，但仍未覆盖物理世界所需的实时状态与行动闭环。

